\section{The Merits Firewall and Administrable Boundaries}
\label{sec:objections}

The front door is an access right with a defined stopping point. It begins only with an adverse event affecting a legislatively specified interest, supplies a closed first answer, and expands only upon a claim-specific showing. It changes no element of the underlying cause of action. Five boundaries make that allocation durable: existing law defines the protected interest; constitutional and confidentiality rules shape disclosure; arbitration, territorial limits, and federal preemption govern the enforcing forum; adoption and deterrence are conceded and priced; and product or speech classification remains for the merits.

\subsection{Existing Legal Interests Define the Gate}

The designation selects an existing legal interest that justifies compulsory observability---for example, employment, credit, housing, public benefits, health, education, physical safety, or the welfare of minors. Stage One requires a verified adverse event affecting that interest, not a complete underlying cause of action or a freestanding demand to audit a business. A rejected applicant identifies the job and disposition; a benefits recipient the entitlement and challenged action; a companion-service user the covered interaction. Capability, dissatisfaction, or a preference for another model does not qualify.

The gate preserves the heterogeneity of AI disputes. Some losses have no private remedy; others sound in contract, tort, antidiscrimination, administrative, consumer-protection, or no-liability rules.\lrfootnote{See \artcite[483, 515--17]{hendersonRobotLiability2026}.} An out-of-scope interest or field ends the access action, as does a legal bar dispositive of the statutory access claim itself, such as lack of jurisdiction, immunity, or express preemption. The absence of a private merits remedy or a limitations bar on an underlying claim does not necessarily foreclose another use permitted by the Act, such as correction, contest, or appeal. Once the first answer identifies the deployment and custodian, ordinary procedure governs. Routine observability and the first answer remain operator costs; extraordinary protocols follow the designation and forum law, and actors outside the control lane exit.

\subsection{Constitutional, Trade-Secret, and Privacy Protections Shape Access}

A compelled answer must satisfy the Constitution as well as the model statute. The first answer is transaction linked and principally factual: what system and version ran, which enterprise holds the record, what output or notice occurred, and which role possessed listed authority. It is delivered to the affected person or enforcing tribunal rather than published to the world. Those features distinguish the first answer from compelled ideological speech or a requirement that an enterprise confess legal fault.

The governing First Amendment standard depends on the speaker, content, and regulatory setting. \emph{Zauderer} permits certain compelled disclosures of purely factual and uncontroversial commercial information reasonably related to preventing deception; \emph{NIFLA} rejects an unlimited category of ``professional speech'' and requires courts to examine burdens and fit.\lrfootnote{See \casecite[650--53]{zauderer1985}; \casecite[768--79]{nifla2018}.} A designation should therefore require operational facts, use neutral field definitions, and avoid compelled normative characterizations such as ``unsafe,'' ``biased,'' or ``responsible.'' If a sector or requested field falls outside the commercial-disclosure line, the otherwise applicable standard of scrutiny governs. The statute's closed fields and nonpublic protections narrow the burden; they do not predetermine the constitutional result.

Trade secrets receive protection, not categorical exclusion. Trade-secret interests can constitute property and turn on reasonable expectations of confidentiality, while litigation law has long used protective orders to limit disclosure of confidential commercial material.\lrfootnote{See \casecite[1001--14]{ruckelshaus1984}; \statcite{frcp26cTradeSecret}.} Stage One does not require source code, weights, or unrestricted training data. It asks for identity, version, event, custodian, and control fields. Those deeper materials may be considered only at Stage Two upon necessity and proportionality, with redaction, attorneys'-eyes-only access, secure expert review, in camera inspection, and testing protocols available. A withholding party identifies the record, custodian, date, version, and legal basis so that protection does not become invisibility.

Privilege follows the same line developed in Section~\ref{sec:mobley-privilege}. Legal advice and attorney work product remain protected. The operator produces underlying nonprivileged facts and ordinary-course business records and logs protected material with sufficient specificity for review. The Act creates no privilege exception and no rule that counsel's involvement automatically removes an operational fact from the minimum record.

Privacy and security limit both creation and production. Part IV's separate clocks and event schema keep unrelated users, raw conversations, worker surveillance, protected health information, credentials, and security details outside the control sheet by default. A request uses the narrowest available event identifier, time window, feature, and population. Pseudonymization, aggregation, secure review, and in camera testing can establish that a control existed without exposing the control itself. A timely field-specific objection pauses only the contested protected field and permits prompt neutral review before disclosure or penalty; uncontested identity and event fields remain due.\lrfootnote{Cf. \casecite[418--23]{patel2015} (requiring an opportunity for neutral precompliance review before penalized administrative inspection). A private first-answer request is not the hotel-inspection regime at issue in \emph{Patel}; the case supplies a protective design principle rather than the governing standard.} The boundary is affirmative: define the factual minimum, minimize collection, and protect what remains.

\subsection{Arbitration, Territorial Limits, and Federal Preemption Govern the Forum}

Nonwaiver describes the substantive duty, not an exclusive judicial forum. An operator cannot contract away the affected person's right to the minimum record. But a valid arbitration agreement may govern who enforces that right. For a federal statutory right, a contrary congressional command can displace the Federal Arbitration Act's rule of arbitrability. A state enactment cannot single out arbitration agreements for disfavored treatment; the Model Act therefore makes the duty nonwaivable while allowing an otherwise valid agreement to select the private enforcement forum.\lrfootnote{See \casecite[98--104]{compuCredit2012}; \casecite[251--58]{kindred2017}; \casecite[359--63]{preston2008}.} Questions about whether the parties delegated arbitrability follow federal arbitration doctrine.\lrfootnote{See \casecite[943--47]{firstOptions1995}.}

The duty still arises before the merits forum needs discovery. A visible operator must maintain the deployment record locally and answer a qualifying request on the statutory timetable; completion is not a condition precedent to commencing or continuing arbitration. If it disputes enforceability or scope, the proper court or arbitrator resolves that dispute under the agreement and the Federal Arbitration Act. An agreement to which the administering agency is not a party does not eliminate public enforcement brought in the agency's own name.\lrfootnote{See \casecite[288--96]{waffleHouse2002}.} This allocation preserves the substantive right, the selected private forum, and independent public enforcement.

The locally maintained deployment record is especially important because arbitral process may not reach an upstream stranger. Federal Arbitration Act section 7 authorizes arbitrators to summon a person to attend and bring documents; the Ninth Circuit reads it not to authorize a stand-alone prehearing document subpoena to a nonparty.\lrfootnote{See \statcite{faaArbitrationSection7}; \casecite[706--09]{cvsVividus2017}.} The operator's predeployment retrieval right supplies the defined vendor fields without enlarging arbitral subpoena power. A supplier directly bound by the Act or contractually bound to respond answers on that source of authority; another supplier responds only through an independent source of law.

Territorial limits remain equally concrete. The State's authority to apply its substantive law, a tribunal's personal jurisdiction, and valid service are separate predicates. Operator notice supplies none of them. A control holder receives a direct statutory duty only when the enterprise and identified deployment fall within valid territorial and choice-of-law reach; coercive relief additionally requires service and an independent jurisdictional basis.\lrfootnote{See \casecite[262--68]{bristolMyers2017}; \casecite[358--72]{fordMotor2021}.} Appointment of the operator's response agent authorizes the statutory delivery route only. It should not silently become consent to all-purpose jurisdiction; an express registration-based consent regime raises the distinct questions illustrated by \emph{Mallory}.\lrfootnote{See \casecite[134--36]{mallory2023}.} The domestic visible operator nevertheless owes the locally maintained control sheet, the event receipt, and the contractual retrieval path. That mandatory local architecture keeps the answer reachable without pretending that state process reaches every global node.

Sectoral designation must also perform a substantive federal-preemption check provision by provision. FCRA expressly preempts specified state requirements concerning dispute procedures, adverse-action duties, and furnisher responsibilities; its present agency interpretation replaced the Bureau's withdrawn 2022 narrow-preemption view, while recognizing that courts remain the final arbiters.\lrfootnote{See \statcite{fcraPreemption}; \statcite{cfpbFCRAPreemption2025}. The agency document is nonbinding guidance; the statutory text controls.} ERISA can displace state reporting regimes that intrude on nationally uniform plan administration. Federal medical-device law preempts some requirements different from or additional to federal requirements, while the HIPAA rules preserve specified more-stringent state privacy provisions and forbid contrary disclosure commands.\lrfootnote{See \statcite{erisaPreemption}; \casecite[320--26]{gobeille2016}; \statcite{medicalDevicePreemption}; \casecite[321--30]{riegel2008}; \statcite{hipaaPreemption}; \statcite{hipaaStateLawPreemption}.} Employment law supplies a different savings rule.\lrfootnote{See \statcite{titleVIISavings}.} A designation should identify federal actors and exclusive remedies, whether federal law supplies a floor, ceiling, or parallel-duty path, which actors and fields are excluded, how remedies overlap, and which claims the State may toll. A generic savings clause cannot preserve an application federal law forecloses. Congress can use the same front-door architecture when national uniformity is required.

\subsection{Two Conceded Trades: Designation Dependency and Thin Deterrence}
\label{sec:trades}

Two limits of the front door are better stated than discovered. First, the machinery is inert until a legislature designates. California's companion-chatbot provisions demonstrate the form without containing the front-door rule, as Part II.A notes; the Article shows that legislatures \emph{can} designate, not that they will designate the fields that bind. The claim here is that the fields are specifiable, administrable, and bounded, not that adoption is assured. Second, the deterrence is thin. Model Act \S~9 provides that bare nonperformance creates no private statutory damages, and Model Act \S~8's merits firewall closes the inference routes that might substitute for damages. A duty that cannot be converted into damages is cheap to ignore; the design buys adoptability---no industry can say the Act creates liability---at the price of deterrent force. And where governing law supplies no private remedy at all, as Part V.A recognizes some losses do, nothing catches an enterprise that exits through the strategic incompetence Part II.E describes. Those are the design's real boundaries, and they are priced here rather than left for a hostile reader to find.

\subsection{Classification Neutrality Preserves the Merits}

The front door does not classify an AI system as a product, service, speaker, agent, or legal person. A companion application can present design questions about access, memory, engagement objectives, crisis response, and warnings while a claim directed at an adopted message presents a distinct First Amendment question. An employment score can participate in a decision without becoming a physical product. Generated text can communicate an operator's decision without making the model a juridical actor.

\emph{Garcia} illustrates the separation. At Rule 12, the district court treated Character.AI as plausibly alleged to be a product because the challenged features concerned design and commercial distribution. It declined on that record to classify the generated sequence as protected speech where defendants had not identified a relevant expressive choice or speaker.\lrfootnote{See \casecite[1174--80]{garcia2025}.} Those rulings were procedural and claim specific. The first answer would identify the deployed system, event, and organizational choices reflected in the bounded record; the pleadings and requested relief identify the challenged conduct and whether the remedy regulates design, transaction, decision process, or expression. Governing law then classifies the claim.

Fragmented and open architectures change the custodian map, not this boundary. An open-weights publisher with no continuing visibility, contractual right, or stopping power can exit after provenance; a modifier that selects the protected use, adds tools, removes safeguards, and controls operation can occupy the lane; and a provider retaining mandatory updates, monitoring, or suspension may retain another. Contract labels and brand association matter only insofar as they describe practical authority. The resulting record lets governing law distinguish reasonable deployment, independent modification, residual risk, and actionable conduct without choosing a winner in advance.
